\cleardoublepage
\phantomsection
\pdfbookmark{Sommario}{Sommario}
\begingroup
\let\clearpage\relax
\let\cleardoublepage\relax
\let\cleardoublepage\relax

\chapter*{Sommario}

Il presente documento descrive il lavoro svolto durante il periodo di stage presso VarGroup di Padova (UAN COMPANY S.r.l), della durata di circa trecento ore nei mesi di maggio e giugno 2026. Gli obiettivi da raggiungere sono riportati di seguito.\\

\noindent L'obiettivo era lo sviluppo di \textit{AI Document Consistency Platform}: una \emph{web application} per l'analisi automatizzata della qualità documentale. La piattaforma consente agli utenti di importare documenti in diversi formati (PDF, Word, Markdown) e di sottoporli a una serie di verifiche, al cui termine viene restituito un \emph{report} dettagliato con le criticità riscontrate e i relativi suggerimenti di correzione. La \emph{web application} prevede inoltre una \emph{dashboard} per la consultazione dello storico dei \emph{report} prodotti.\\
Il processo di analisi si articola in un sistema ibrido a due livelli:
\begin{itemize}
    \item Il primo livello, di natura deterministica, esegue controlli algoritmici sul documento basandosi su un insieme di regole configurabili (\emph{ruleset}). In questo passaggio vengono validati i limiti di lunghezza imposti per le sezioni;
    \item Il secondo livello utilizza modelli \emph{\gls{llm}}\glsfirstoccur, nello specifico i modelli Claude di Anthropic, con l'obiettivo di individuare incoerenze logiche, contraddizioni tra paragrafi, violazioni delle \emph{policy} aziendali e assenza delle sezioni obbligatorie.
\end{itemize}

\noindent Infine, è stata richiesta la produzione della documentazione tecnica, composta dalle seguenti parti:
\begin{itemize}
    \item \emph{README} per la configurazione e l'avvio del progetto in 
    ambiente locale e tramite \emph{\gls{docker}}\glsfirstoccur;
    \item Documentazione delle \emph{\gls{api}}\glsfirstoccur tramite \emph{\gls{swagger}}\glsfirstoccur;
    \item Guida alla configurazione dei \emph{ruleset}.
\end{itemize}
 
%\vfill

%\selectlanguage{english}
%\pdfbookmark{Abstract}{Abstract}
%\chapter*{Abstract}

%\selectlanguage{italian}

\endgroup

\vfill
